% ECONOMICS_PROBLEM: {"id": "D91-507", "title": "Equilibrium effects of nudges at scale", "jel_code": "D91", "source_id": "OP-0066"}
% FIRST_POSTED: 2026-10-09
% ATTEMPT_STATUS: partial
% ENTRY_KIND: research_attempt
\documentclass[11pt,a4paper]{article}
\usepackage[T1]{fontenc}
\usepackage[utf8]{inputenc}
\usepackage[margin=27mm]{geometry}
\usepackage{amsmath,amssymb,amsthm,mathtools}
\usepackage[hidelinks]{hyperref}
\setlength{\parindent}{0pt}
\setlength{\parskip}{5pt}
\newtheorem{theorem}{Theorem}[section]
\newtheorem{proposition}[theorem]{Proposition}
\newtheorem{lemma}[theorem]{Lemma}
\newcommand{\R}{\mathbb R}
\newcommand{\E}{\mathbb E}
\DeclareMathOperator{\Var}{Var}
\DeclareMathOperator{\Cov}{Cov}
\begin{document}
\section*{Nudges at scale: equilibrium attenuation, amplification, and welfare}
\section{Scope}
An individual experiment at one market saturation identifies a policy embedded in that equilibrium. It does not automatically identify full-coverage effects. This attempt solves two explicit market benchmarks, gives sharp extrapolation bounds under a stated smoothness restriction, and derives finite-sample simultaneous inference for randomized saturations. The general market game, endogenous entry and evaluator-dependent welfare in the catalogue remain unresolved.

First consider a unit mass of consumers with quantity demands
\[
 q_i=a-bp+\tau z_i+\gamma\overline q,\qquad b>0,
\]
where $z_i\in\{0,1\}$ is assignment and $s=E[z_i]$. Competitive supply is $p=c+\kappa\overline q$, $\kappa\ge0$. All compared quantities are assumed nonnegative and interior; these linear equations are not extrapolated beyond that domain. Market aggregation and substitution give
\[
 \overline q(s)=\frac{a-bc+\tau s}{1-\gamma+b\kappa}. \tag{1}
\]
Assume $D=1-\gamma+b\kappa>0$. The algebraic equilibrium is then unique. If an adaptive best-response process is intended, its convergence requires its own stability condition; uniqueness alone does not prove dynamical stability.

\section{An exact comparison between direct and rollout effects}
Holding the market state fixed, the individual assignment contrast equals $\tau$. In the continuum model one consumer has no effect on the aggregate, so this is also the direct contrast within a fixed-saturation equilibrium. The full market change is instead
\[
 \overline q(1)-\overline q(0)=\tau/D. \tag{2}
\]
Thus supply feedback attenuates the effect when $D>1$, while positive peer feedback can amplify it when $0<D<1$. With $b\kappa=1,\gamma=0$, rollout equals half the direct response. With $b\kappa=0.2,\gamma=0.8$, it equals $2.5$ times the direct response. These are feasible constructed coefficients, subject to choosing $a,c,\tau$ so quantities remain interior.

For a treated consumer the total change from the zero-coverage state includes both its own $\tau$ and the common price/peer term. Untreated consumers also change quantities through that common term. Therefore a difference between treated and untreated consumers can remove precisely the equilibrium effect relevant to the market mean. Randomizing coverage across independently operating markets supplies the variation needed to observe (1) as a function of $s$; randomizing only individuals at one $s$ does not.

\section{More quantity need not mean more welfare}
For a separate homogeneous benchmark, let the true marginal willingness to pay be $a-bQ$ and real marginal supply cost $c+\kappa Q$, with $b>0,\kappa\ge0$. Consumers perceive an intercept $a-u+\tau s$, where $u\ge0$ is a baseline underestimation and the nudge corrects perceived value by $\tau s$. Then
\[
 Q(s)=\frac{a-c-u+\tau s}{b+\kappa},\qquad
 Q^*=\frac{a-c}{b+\kappa}.
\]
Require interior quantities over the examined range. The evaluator uses true utility, includes producer surplus, and subtracts real implementation cost $ks$, $k\ge0$:
\[
 W(s)=(a-c)Q(s)-\tfrac12(b+\kappa)Q(s)^2-ks.
\]
Payments cancel between consumers and producers. Completing the square yields
\[
 W(s)-W(0)=\frac{u\tau s-\tfrac12\tau^2s^2}{b+\kappa}-ks. \tag{3}
\]
This is an exact welfare comparison, not an assertion that the chosen evaluator is uniquely correct. With $u=0$, any nonzero demand-raising nudge reduces surplus before costs because the benchmark started at the efficient quantity. With $u>0$, initial correction may help but excessive coverage can overcorrect. If $\tau>0$, the welfare-maximizing coverage on $[0,1]$ is
\[
 s^*=\operatorname{clip}_{[0,1]}
 \left(\frac{u}{\tau}-\frac{k(b+\kappa)}{\tau^2}\right).
\]
Different privacy, inconvenience or moral-pressure costs alter (3) and the optimum. The first quantity model's peer term cannot be inserted into this welfare formula without defining whether peer effects enter true utility or merely perceived utility; these are intentionally separate benchmarks.

\section{Sharp bounds when full saturation was not observed}
Let $f(s)$ be the population market-level outcome under the declared equilibrium selection rule. Suppose experiments identify $f(s_j)=y_j$ at finitely many saturation levels and economic reasoning supplies a global Lipschitz constant $L$ on $[0,1]$. Assume the observations satisfy $|y_i-y_j|\le L|s_i-s_j|$. Then
\[
 f(1)\in\left[
  \max_j\{y_j-L(1-s_j)\},\quad
  \min_j\{y_j+L(1-s_j)\}\right]. \tag{4}
\]
\begin{proposition}
The bounds in (4) are sharp over all real-valued $L$-Lipschitz functions matching the experimental means.
\end{proposition}
\begin{proof}
Necessity follows from the Lipschitz inequalities between $1$ and every observed level. For any value $v$ inside (4), augment the finite data by $(1,v)$. All pairwise inequalities still hold. The function
$F(s)=\min_i\{y_i+L|s-s_i|\}$, with the augmented point included, is $L$-Lipschitz and agrees with every augmented value. Each cone has the same Lipschitz constant, and the pairwise inequalities force equality at the data points. This constructs an admissible extension attaining $v$.
\end{proof}
Known outcome support can be enforced by clipping the extension to that support, preserving both the Lipschitz constant and values already in the support. An untested assumed $L$ is not identified by observing a few small saturations. With no continuity restriction, $f(1)$ can take any supported value independently of the observed interior levels. Multiple equilibria may also make a selected $f$ discontinuous; a Lipschitz assumption on primitives does not automatically bound a selection rule that jumps.

\section{Finite-sample inference at the market level}
Suppose each tested saturation has $n_j$ independent markets whose bounded outcomes belong to $[A,B]$, and the assignment design identifies its target mean. With $J$ saturation levels,
\[
 r_j=(B-A)\sqrt{\log(2J/\alpha)/(2n_j)}
\]
gives simultaneous intervals $[\widehat y_j-r_j,\widehat y_j+r_j]$ with coverage at least $1-\alpha$, by Hoeffding and a union bound. Empty cells use the full outcome range. Intersect the set of all Lipschitz functions with these interval constraints and project to $f(1)-f(0)$. On the common event, the true function belongs to the feasible set, hence the projection covers the rollout contrast. This projection also preserves the coupling between baseline and full saturation. Separate coordinate confidence limits can be combined conservatively, but should not be called a sharp joint confidence region.

The procedure remains valid near zero direct effects because it uses no division by an estimated response. It can be wide if all observed saturations are low. Larger cluster counts improve sampling precision but cannot remove the distance-to-full-coverage term in (4) without stronger restrictions or high-saturation data.

\section{Unresolved empirical and dynamic content}
The results show precisely how a supply slope, peer effect or welfare evaluator changes an extrapolation. A completed catalogue application needs markets, an operational nudge, equilibrium adjustment time, supplier quality and entry data, and validation away from fitted saturations. Individual take-up effects, market quantity effects and social welfare are separate outcomes. None of the formulas claims a universal attenuation factor for nudges.
\begin{thebibliography}{9}
\bibitem{ak} H. Allcott and J. B. Kessler (2019), \emph{The Welfare Effects of Nudges: A Case Study of Energy Use Social Comparisons}, AEJ: Applied Economics 11(1), 236--276. \url{https://www.aeaweb.org/articles?id=10.1257/app.20170328}.
\bibitem{dl} S. DellaVigna and E. Linos (2022), \emph{RCTs to Scale: Comprehensive Evidence from Two Nudge Units}, Econometrica 90, 81--116. \url{https://doi.org/10.3982/ECTA18709}. The cited scale comparison does not supply the equilibrium coefficients assumed above.
\end{thebibliography}

\end{document}
